\section{Пакеты прикладных программ для математического моделирования}

Пакеты прикладных программ являются одним из основных инструментов вычислительного
математического моделирования. Они позволяют перейти от математической постановки
задачи к её численному решению, анализу результатов и визуализации без необходимости
реализовывать весь вычислительный алгоритм с нуля.

При этом использование готового пакета не устраняет необходимость понимать
математическую модель и численный метод. Пользователь должен уметь оценивать
применимость выбранной модели, корректность постановки задачи, свойства численного
алгоритма и достоверность полученного решения.


\subsection{Понятие пакета прикладных программ}

\textbf{Пакет прикладных программ (ППП)} --- совокупность программных средств,
предназначенных для решения определённого класса прикладных задач и объединённых
общими средствами задания исходных данных, вычислений и обработки результатов.

В математическом моделировании ППП обычно обеспечивает некоторую или всю цепочку
вычислительного эксперимента:
\[
\text{математическая модель}
\longrightarrow
\text{дискретная модель}
\longrightarrow
\text{численный алгоритм}
\longrightarrow
\text{программная реализация}
\longrightarrow
\text{результат}.
\]

В отличие от отдельной программы, пакет обычно содержит:

\begin{itemize}
	\item набор готовых вычислительных алгоритмов;
	\item средства задания параметров модели;
	\item средства работы с геометрией и расчётными сетками;
	\item библиотеки математических функций;
	\item средства визуализации;
	\item средства обмена данными;
	\item интерфейсы для расширения и автоматизации расчётов.
\end{itemize}

Примерами различных классов таких систем являются математические среды
MATLAB, Mathematica, Maple, системы инженерного моделирования ANSYS и COMSOL,
а также специализированные открытые комплексы, например OpenFOAM.

Следует различать \textbf{математическую модель}, \textbf{численный метод}
и \textbf{программный пакет}. Пакет является инструментом реализации модели,
но сам по себе не гарантирует правильность постановки или физическую адекватность
полученного результата.


\subsection{Классификация пакетов математического моделирования}

Пакеты можно классифицировать по нескольким признакам.

\textbf{По назначению:}

\begin{enumerate}
	\item \textbf{Универсальные математические среды.}
	
	Содержат средства линейной алгебры, решения ОДУ, оптимизации, обработки данных,
	статистики и визуализации.
	
	\item \textbf{Системы компьютерной алгебры.}
	
	Ориентированы прежде всего на символьные преобразования:
	дифференцирование, интегрирование, решение уравнений, преобразование выражений,
	работу с рядами и специальными функциями.
	
	\item \textbf{CAE-системы} (\textit{Computer-Aided Engineering}).
	
	Предназначены для численного моделирования инженерных объектов:
	механики твёрдого тела, теплопроводности, гидродинамики,
	электромагнитных задач и т.д.
	
	\item \textbf{Специализированные вычислительные пакеты.}
	
	Предназначены для конкретного класса задач, например вычислительной
	гидродинамики, молекулярной динамики, оптимизации или обработки экспериментальных
	данных.
	
	\item \textbf{Библиотеки численных методов.}
	
	Не обязательно имеют графический интерфейс, а предоставляют программные
	реализации алгоритмов для использования в пользовательских программах.
\end{enumerate}

\textbf{По способу взаимодействия} можно выделить:

\begin{itemize}
	\item пакеты с графическим пользовательским интерфейсом;
	\item пакеты с командным интерфейсом;
	\item программные библиотеки;
	\item гибридные системы, совмещающие GUI, язык сценариев и API.
\end{itemize}

\textbf{По степени открытости} различают коммерческие и свободно распространяемые
пакеты, а также системы с открытым и закрытым исходным кодом.


\subsection{Архитектура и основные компоненты}

Современный программный комплекс математического моделирования обычно имеет
модульную архитектуру.

Основные компоненты можно представить в виде цепочки
\[
\text{препроцессор}
\rightarrow
\text{решатель}
\rightarrow
\text{постпроцессор}.
\]

\textbf{Препроцессор} используется для подготовки задачи. В нём выполняются:

\begin{itemize}
	\item построение или импорт геометрии;
	\item задание физических свойств;
	\item выбор математической модели;
	\item задание начальных и граничных условий;
	\item построение вычислительной сетки;
	\item настройка параметров численного метода.
\end{itemize}

\textbf{Решатель} является вычислительным ядром системы. Он формирует и решает
дискретную задачу.

Например, после дискретизации дифференциального уравнения может возникнуть система
линейных алгебраических уравнений
\[
A x = b,
\]
либо нелинейная система
\[
F(x)=0.
\]

Для нестационарной задачи решатель последовательно вычисляет приближения
\[
u^{0},u^{1},\ldots,u^{N},
\]
соответствующие различным моментам времени.

\textbf{Постпроцессор} предназначен для анализа результатов:

\begin{itemize}
	\item построения графиков;
	\item визуализации полей;
	\item вычисления интегральных характеристик;
	\item анализа экстремальных значений;
	\item сравнения нескольких расчётов;
	\item экспорта результатов.
\end{itemize}

Кроме этих трёх компонентов в пакет могут входить:

\begin{itemize}
	\item база материалов и физических моделей;
	\item генератор расчётных сеток;
	\item модуль оптимизации;
	\item средства параллельных вычислений;
	\item язык сценариев;
	\item программный интерфейс API;
	\item система управления расчётными проектами.
\end{itemize}


\subsection{Численные решатели}

Центральной частью пакета математического моделирования являются
\textbf{численные решатели}.

В зависимости от задачи могут использоваться различные классы алгоритмов.

\textbf{Решение систем линейных уравнений:}

\begin{itemize}
	\item метод Гаусса;
	\item $LU$-разложение;
	\item метод Холецкого;
	\item итерационные методы;
	\item методы подпространств Крылова.
\end{itemize}

Для больших разреженных систем особенно важны итерационные методы и
предобуславливатели.

\textbf{Решение нелинейных систем}

\[
F(x)=0
\]

может выполняться методом Ньютона и его модификациями:
\[
J(x^{(k)})\Delta x^{(k)}
=
-F(x^{(k)}),
\]
\[
x^{(k+1)}
=
x^{(k)}+\Delta x^{(k)},
\]
где $J$ --- матрица Якоби.

\textbf{Для обыкновенных дифференциальных уравнений}

\[
\frac{dy}{dt}=f(t,y)
\]

используются, например,

\begin{itemize}
	\item методы Эйлера;
	\item методы Рунге--Кутты;
	\item методы Адамса;
	\item BDF-методы.
\end{itemize}

\textbf{Для уравнений в частных производных} применяются:

\begin{itemize}
	\item метод конечных разностей;
	\item метод конечных элементов;
	\item метод конечных объёмов;
	\item метод граничных элементов;
	\item спектральные методы.
\end{itemize}

Важными характеристиками решателя являются:

\begin{itemize}
	\item порядок аппроксимации;
	\item устойчивость;
	\item скорость сходимости;
	\item вычислительная сложность;
	\item требования к памяти;
	\item возможность параллельного выполнения.
\end{itemize}

Пользователь должен контролировать критерии остановки итерационного процесса.
Например, решение может считаться найденным при выполнении условия
\[
\|r^{(k)}\| < \varepsilon,
\]
где
\[
r^{(k)}=b-Ax^{(k)}
\]
--- невязка, а $\varepsilon$ --- заданная точность.

Таким образом, получение программой некоторого числа ещё не означает
получения корректного решения: необходимо контролировать сходимость,
погрешность и устойчивость вычислительного процесса.


\subsection{Подготовка исходных данных и постобработка}

Качество результата в значительной степени определяется качеством исходных данных.

При подготовке задачи задаются:

\begin{itemize}
	\item геометрические параметры;
	\item физические свойства;
	\item коэффициенты математической модели;
	\item начальные условия;
	\item граничные условия;
	\item параметры дискретизации;
	\item параметры численного решателя.
\end{itemize}

Для пространственных задач важным этапом является построение
\textbf{вычислительной сетки}.

Сетка должна быть достаточно подробной в областях, где решение быстро изменяется,
например:

\begin{itemize}
	\item вблизи границ;
	\item около концентраторов напряжений;
	\item в пограничных слоях;
	\item в областях сильных градиентов температуры или скорости.
\end{itemize}

Однако чрезмерное сгущение сетки увеличивает вычислительные затраты.

Поэтому выполняют исследование \textbf{сеточной сходимости}:
решают одну и ту же задачу на последовательности сеток
\[
h_1 > h_2 > h_3 > \ldots
\]
и проверяют стабилизацию интересующей величины:
\[
Q_{h_1},Q_{h_2},Q_{h_3},\ldots
\rightarrow Q.
\]

После расчёта проводится постобработка.

Для скалярного поля
\[
u=u(x,y,z)
\]
можно строить линии или поверхности уровня
\[
u(x,y,z)=C.
\]

Для векторных полей, например скорости
\[
\mathbf{v}=(v_x,v_y,v_z),
\]
используются стрелочные диаграммы, линии тока и распределение модуля
\[
|\mathbf v|
=
\sqrt{v_x^2+v_y^2+v_z^2}.
\]

При постобработке важно анализировать не только красивую визуализацию,
но и количественные характеристики: интегралы, нормы ошибок, балансы,
экстремумы и зависимости от параметров.


\subsection{Форматы данных и интеграция}

В реальном вычислительном проекте один программный пакет редко используется
изолированно.

Необходим обмен:

\begin{itemize}
	\item геометрическими моделями;
	\item расчётными сетками;
	\item табличными данными;
	\item матрицами и векторами;
	\item полями физических величин;
	\item результатами вычислительного эксперимента.
\end{itemize}

Форматы можно разделить на:

\begin{itemize}
	\item \textbf{текстовые} --- CSV, JSON, XML и другие;
	\item \textbf{бинарные} --- эффективны для больших массивов данных;
	\item \textbf{специализированные научные форматы};
	\item \textbf{форматы геометрии и расчётных сеток}.
\end{itemize}

Текстовые форматы просты для анализа и переноса, но при больших объёмах данных
занимают много места и медленнее обрабатываются.

Бинарные форматы обеспечивают более компактное хранение и быстрый ввод-вывод,
но требуют согласованного программного интерфейса.

Современные пакеты часто предоставляют \textbf{API}
(\textit{Application Programming Interface}), позволяющий управлять расчётом
из внешней программы.

Например, можно автоматизировать серию вычислительных экспериментов:
\[
p_1,p_2,\ldots,p_N
\longrightarrow
Q(p_1),Q(p_2),\ldots,Q(p_N),
\]
где $p_i$ --- различные значения параметра модели.

Это особенно важно при:

\begin{itemize}
	\item оптимизации;
	\item идентификации параметров;
	\item анализе чувствительности;
	\item обработке большого числа вариантов;
	\item построении суррогатных моделей.
\end{itemize}


\subsection{Верификация программного комплекса}

Наличие готового программного пакета не освобождает пользователя от проверки
правильности вычислений.

Следует различать два понятия:

\begin{itemize}
	\item \textbf{верификация} --- проверка того, что математическая и численная
	задача реализована программой правильно;
	\item \textbf{валидация} --- проверка того, что сама математическая модель
	достаточно хорошо описывает реальный объект.
\end{itemize}

Условно:

\[
\text{верификация: «правильно ли решаем уравнения?»},
\]

\[
\text{валидация: «правильные ли уравнения решаем?»}.
\]

Для верификации программного комплекса применяются следующие подходы.

\textbf{1. Сравнение с аналитическим решением.}

Если для тестовой задачи известно точное решение $u(x)$, вычисляют погрешность
\[
e_h = u_h-u
\]
и её норму
\[
\|e_h\|.
\]

\textbf{2. Исследование сходимости.}

При уменьшении шага дискретизации должно наблюдаться
\[
\|u-u_h\|\rightarrow 0,
\qquad h\rightarrow 0.
\]

Если метод имеет порядок $p$, ожидается поведение
\[
\|u-u_h\|\approx C h^p.
\]

\textbf{3. Метод изготовленных решений.}

Заранее выбирается функция $u(x)$, после чего по ней вычисляется такая правая часть
уравнения, чтобы выбранная функция являлась его точным решением.
Затем проверяется способность программы воспроизвести это решение с ожидаемой
погрешностью.

\textbf{4. Сравнение с независимой программной реализацией.}

Одна задача решается различными программами или различными алгоритмами.

\textbf{5. Проверка законов сохранения.}

Например, для замкнутой физической системы контролируются баланс массы,
импульса или энергии.

\textbf{6. Регрессионное тестирование.}

После изменения программного кода повторяются заранее выбранные тестовые задачи
и проверяется, что результаты не изменились неожиданным образом.

Кроме того, следует контролировать:

\begin{itemize}
	\item независимость результата от сетки;
	\item влияние временного шага;
	\item величину невязки;
	\item чувствительность к параметрам решателя;
	\item корректность единиц измерения;
	\item выполнение начальных и граничных условий.
\end{itemize}


\subsection{Критерии выбора пакета}

Выбор программного пакета должен определяться прежде всего задачей, а не
популярностью конкретного продукта.

Основными критериями являются:

\begin{enumerate}
	\item \textbf{Соответствие математической модели.}
	
	Пакет должен поддерживать необходимые уравнения, физические модели
	и типы граничных условий.
	
	\item \textbf{Наличие подходящих численных методов.}
	
	Важно учитывать устойчивость, точность и применимость решателя
	к рассматриваемой задаче.
	
	\item \textbf{Размер задачи.}
	
	Для больших задач существенны работа с разреженными матрицами,
	параллельные вычисления и экономное использование памяти.
	
	\item \textbf{Возможность работы на многопроцессорных системах.}
	
	Для ресурсоёмких моделей важна масштабируемость на CPU и GPU.
	
	\item \textbf{Средства построения сеток и постобработки.}
	
	Для инженерных задач эти возможности могут быть не менее важны,
	чем сам численный решатель.
	
	\item \textbf{Расширяемость.}
	
	Желательна возможность добавления собственных моделей,
	алгоритмов и пользовательских функций.
	
	\item \textbf{Интеграция с другими программами.}
	
	Важны поддерживаемые форматы данных, API и средства автоматизации.
	
	\item \textbf{Воспроизводимость расчётов.}
	
	Необходимо иметь возможность сохранить параметры модели,
	сетку, настройки решателя и версию программного комплекса.
	
	\item \textbf{Лицензирование и стоимость.}
	
	Для научных и промышленных проектов условия лицензии могут существенно
	влиять на выбор системы.
	
	\item \textbf{Документация и поддержка.}
	
	Наличие документации, тестовых задач и активного сообщества
	существенно облегчает проверку и сопровождение модели.
\end{enumerate}

Таким образом, пакет прикладных программ следует рассматривать как
\textbf{инструмент вычислительного эксперимента}. Его основная ценность состоит
не просто в возможности получить численный результат, а в обеспечении полного
цикла:
\[
\boxed{
	\text{постановка задачи}
	\rightarrow
	\text{дискретизация}
	\rightarrow
	\text{решение}
	\rightarrow
	\text{анализ}
	\rightarrow
	\text{проверка результата}
}
\]

При этом ответственность за выбор модели, численного метода и интерпретацию
результатов остаётся за исследователем.

\newpage